\documentclass[10pt,a4paper]{amsart}
\usepackage[margin=19mm]{geometry}
\usepackage{amsmath,amssymb,mathtools}
\usepackage[scr=rsfs,cal=euler]{mathalfa}
\usepackage{xfrac}
\usepackage[unicode,hidelinks,bookmarks=false]{hyperref}
\newcommand{\ie}{i.e.\ }
\newcommand{\cf}{cf.\ }
\newcommand{\resp}{respectively\ }
\newcommand{\Subsection}{\subsection}
\providecommand{\nobreakdash}{\nobreak\hbox{-}}
\setlength{\emergencystretch}{2em}
\multlinegap0pt
\usepackage{amssymb}
\newtheorem{theorem}{Theorem}
\title{weakly $Q$-ideals of commutative rings}
\author{Mahdi Anbarloei}
\date{Source: arXiv:2606.08055v1 (CC BY 4.0)}
\begin{document}
\maketitle

\noindent\emph{Source and licence.} weakly $Q$-ideals of commutative rings. Version arXiv:2606.08055v1 (CC BY 4.0), distributed by arXiv under the Creative Commons Attribution 4.0 International licence, http://creativecommons.org/licenses/by/4.0/, as recorded in the arXiv licence metadata for this exact version and shown on the abstract page https://arxiv.org/abs/2606.08055v1. Full licence notice: \texttt{licenses/arxiv-2606.08055-CC-BY-4.0.txt}.\par\smallskip

\noindent\textit{Editorial scope.} Counted unit: the article's theorem 4 (source0.tex lines 149--158). The article's complete original proof of it is delivered with it (source0.tex lines 159--172, one \texttt{proof} environment of 14 lines). No mathematics is written, repaired or re-derived: the statement and the proof are the article's own lines, in the article's own notation, and no retained line is substituted or re-flowed.  No further source-local context is needed: every cross-reference of every retained line resolves inside the two delivered spans.  Nothing was omitted from any retained span, so the package carries no omission record.  No retained line contains a citation, so the wrapper carries no bibliography and no bibliography entry.  No retained line contains a figure environment, an image-inclusion command or drawing code, so no image is delivered and the package declares no asset dependency.  Editorial formatting only, in addition: an AMS \texttt{amsart} preamble that restates the article's own declarations and macros used by the retained lines, namely the environments \texttt{theorem} on separate counters, together with the standard packages those lines need.  The retained environments print in this document as Theorem 1 in order of appearance, whereas the article prints the same items with the numbering of its own full document.  The complete native source of the article as distributed in the arXiv e-print for this exact version is bundled unchanged as \texttt{sources/202298/source0.tex}.
\begin{theorem}\label{4}
Let $I$ and $Q$ be  ideals  of $A$. Then, the following statements are equivalent:
\begin{enumerate}
\item $I$ is a weakly $Q$-ideal of $A$.
\item $(I:x)=I \cup (0:x)$ for each $x \in A \backslash  Q$.
\item $(I:x) \subseteq Q \cup (0:x)$ for each $x \in A \backslash I$.
\item If $x \in A$ and $I_1$ is an ideal of $A$ with $0 \neq I_1x\subseteq I$, then $I_1 \subseteq Q$ or $x \in I$.
\item If $I_1$ and $I_2$ are ideals of $A$ with $0 \neq I_1I_2 \subseteq I$, then $I_1 \subseteq Q$ or $I_2 \subseteq I$.
\end{enumerate}
\end{theorem}

\begin{proof}
(1) $\Longrightarrow $ (2) Assume that $x \in A \backslash Q$. The inclusion $I \cup (0 : x) \subseteq (I:x)$ holds. Let $y \in (I : x)$. Then, we have $xy \in I$. If $xy \neq 0$, then we get $y \in I$ as $I$ is a weakly $Q$-ideal of $A$. If $xy=0$, then we obtain $y \in (0:x)$ which means $(I:x) \subseteq I \cup (0:x)$. Thus, $(I:x)=I \cup (0:x)$.

(2) $\Longrightarrow $ (1) Assume that $0 \neq xy \in I$ for $x,y \in A$ and $x \notin Q$. By the hypothesis, $y \in I \cup (0 : x)$. Since $y \notin (0:x)$, we have $y \in I$. 

(1) $\Longrightarrow $ (3) By a similar argument to (1) $\Longrightarrow $ (2), we conclude that  the claim is true.

(3) $\Longrightarrow $ (4) Let $0 \neq I_1 x \subseteq I$    for some ideal $I_1$ of $A$  such that $x \in A \backslash I$. Since  $(I : x) \neq (0 : x)$, we obtain  $(I : x) \subseteq Q$ by the hypothesis. Hence, $I_1 \subseteq Q$.


(4) $\Longrightarrow $ (5) Let $0 \neq I_1I_2 \subseteq I$ for some ideals $I_1$ and $I_2$  of $A$ such that neither $I_1 \subseteq Q$ nor $I_2 \subseteq I$. From $0 \neq I_1I_2 \subseteq I$ it follows that $0 \neq I_1 x \subseteq I$ for some $x \in I_2$. This implies that $x \in I$ as $I_1 \nsubseteq Q$. Take any $y \in I_2 \backslash I$. If $I_1 y \neq 0$, then we obtain $y \in I$, a contradiction. Therefore, $I_1 y=0$. Since $0 \neq I_1(x+y) \subseteq I$ and $I_1 \nsubseteq Q$, we obtain $x+y \in I$ and so $y \in I$, a contradiction. 

(5) $\Longrightarrow $ (1) Let $0 \neq xy \in I$ for $x,y \in A$. Set $I_1=\langle x \rangle$ and $I_2=\langle y \rangle$. Then, $0 \neq I_1 I_2 \subseteq I$. By the hypothesis, we conclude that $x \in I_1 \subseteq Q$ or $y \in I_2 \subseteq I$. 
\end{proof}

\end{document}
